%!TEX root=../../main.tex

\begin{lemma}\label{lemma:unique-fit-cgl}
	The model fit is the same for all optimal solutions to the
	CGL problem.
	That is,
	\[
		X w = X w' \quad \text{ for all } w, w' \in \solfn(\lambda).
	\]
	As a consequence, the sum of group norms is also constant
	when \( \lambda > 0 \).
\end{lemma}
\begin{proof}
	This follows in a similar fashion to the classic result for the group lasso.
	Let \( w, w' \in \solfn(\lambda) \), \( \bar w = \half w + \half w' \),
	and suppose \( p^* \) is the optimal value of the constrained program.
	By convexity, we have
	\begin{align*}
		\half \norm{X \bar w - y}_2^2 + \lambda \sum_{\bi \in \calB} \norm{\bwi}_2
		 & \leq
		\half p^* + \half p^* = p^*,
	\end{align*}
	where the inequality is strict if \( X w \neq X w' \) by strong convexity
	of \( f(u) = \norm{u - y}_2^2 \).
	Since \( w, w' \) are both feasible, \( \bar w \) is also feasible and
	clearly \( \bar w \) cannot obtain an objective value less than \( p^* \).
	Thus, \( X w = X w' \)  must hold.

	To see the second part of the result, observe that
	\[
		\lambda \sum_{\bi \in \calB} \norm{\bwi}_2 = p^* - \half \norm{\hat y(\lambda)  - y}_2^2,
	\]
	is also constant over \( \solfn(\lambda) \).
\end{proof}

\solFnCGL*
\begin{proof}
	Fix \( \lambda > 0 \) and let
	\( w \in \solfn(\lambda) \).
	If \( \wi \neq 0 \), then the KKT conditions require
	\[
		\lambda \frac{\wi}{\norm{\wi}_2} = \vi \implies \wi = \alpha_\bi \vi,
	\]
	for \( \alpha > 0 \).
	If \( \wi = 0 \), then \( \wi = \alpha_\bi \vi \) holds trivially for
	\( \alpha_\bi = 0 \).
	Since \( w \) is optimal, it must satisfy
	\[
		X \w = \hat y,
	\]
	by \cref{lemma:unique-fit-cgl}.
	Finally, \( \act(w) \subseteq \calS_\lambda \) so that
	\( w \) satisfies the characterization.

	For the reverse direction, we start by defining
	\[
		\begin{aligned}
			\calX =
			\bigg\{ w \in \R^d : \,
			 & \forall \, \bi \in \calS_\lambda, \,
			\wi = \alpha_\bi \vi, \alpha_\bi \geq 0, \,          \\
			 & \forall \, b_j \in \calB \setminus \calS_\lambda,
			\w_{b_j} = 0, \, X w = \hat y
			\bigg\}
		\end{aligned}
	\]
	Take any \( w' \in \calX \).
	If \( \wi' \neq 0 \), then \( \bi \in \equi \) and \( \norm{\vi}_2 = \lambda \)
	so that we may write
	\begin{align*}
		\wi' = \alpha_\bi \vi
		 & \implies \frac{\wi'}{\norm{\wi'}_2} = \frac{\alpha_\bi \vi}{\norm{\alpha_\bi \vi}_2} \\
		 & \implies \frac{\wi'}{\norm{\wi'}_2} = \frac{\vi}{\lambda}                            \\
		 & \implies \lambda \frac{\wi'}{\norm{\wi'}_2} = \vi.
	\end{align*}
	That is,
	\[
		\Xbi^\top (Xw - y) + \lambda \frac{\wi'}{\norm{\wi'}_2} + \Ki \rmin_\bi = 0,
	\]
	which is exactly stationarity of the Lagrangian.

	If \( \wi' = 0 \), then \( X w' = \hat y \) implies
	\[
		\norm{\Xbi^\top (y - Xw') + K \rmin_\bi}_2
		= \norm{\Xbi^\top (y - Xw) + K \rmin_\bi}_2
		\leq \lambda,
	\]
	for some \( w \in \solfn(\lambda) \),
	which also implies the Lagrangian is stationary.

	Now we show optimality of
	\( w' \) by checking feasibility and complementary slackness.
	If \( \wi' \neq 0  \)
	then \( \wi' = \frac{\alpha_\bi'}{\alpha_\bi} \wi \)
	for some optimal solution \( w \) with \( \alpha_\bi > 0 \).
	This follows since \( \wi' \neq 0 \) implies \( \bi \in \calS_\lambda \).
	Thus,
	\[
		\Ki^\top \wi' = \frac{\alpha_\bi'}{\alpha_\bi} \wi \leq 0,
	\]
	by feasibility of \( \wi \).
	Similarly, we find that complementary slackness is satisfied as follows:
	\[
		[\ri]_j \cdot [\Ki]_j^\top \wi'
		= \frac{\alpha_\bi'}{\alpha_\bi} [\ri]_j \cdot [\Ki]_j^\top \wi
		= 0.
	\]
	If \( \wi = 0 \), then both feasibility and complementary slackness
	are trivial.
	Since \( (w', \rmin) \) are feasible and \( \rmin \) is dual optimal,
	we conclude the KKT conditions are satisfied and thus
	\( w' \in \solfn(\lambda) \).
	This completes the proof.
\end{proof}

\begin{proposition}\label{prop:cgl-alternate}
	Fix \( \lambda > 0 \).
	The optimal set for CGL problem is also
	given by
	\begin{equation}\label{eq:sol-fn-cgl-alternate}
		\begin{aligned}
			\solfn(\lambda) =
			\bigg\{ w \in \R^d : \,
			 & \forall \, \bi \in \equi, \,
			\wi = \alpha_\bi \vi, \alpha_\bi \geq 0, \,                   \\
			 & \forall \, b_j \in \calB \setminus \equi, \w_{b_j} = 0, \,
			X w = \hat y,                                                 \\
			 & \, K^\top w \leq 0, \,
			\abr{\rmin, K^\top w} = 0
			\bigg\}
		\end{aligned}
	\end{equation}
\end{proposition}
\begin{proof}
	Fix \( \lambda > 0 \) and let
	\( w \in \solfn(\lambda) \).
	If \( \wi \neq 0 \), then the KKT conditions require
	\[
		\lambda \frac{\wi}{\norm{\wi}_2} = \vi \implies \wi = \alpha_\bi \vi,
	\]
	for \( \alpha > 0 \).
	If \( \wi = 0 \), then \( \wi = \alpha_\bi \vi \) holds trivially for
	\( \alpha_\bi = 0 \).
	Since \( w \) is optimal, it must satisfy
	\[
		X \w = \hat y,
	\]
	by \cref{lemma:unique-fit-cgl}.
	Finally, \( \abr{\rmin, K w} = 0 \) and is feasible by KKT conditions
	so that \( w \) satisfies the characterization.

	For the reverse direction, we start by defining
	\[
		\begin{aligned}
			\calX =
			\bigg\{ w \in \R^d : \,
			 & \forall \, \bi \in \equi, \,
			\wi = \alpha_\bi \vi, \alpha_\bi \geq 0, \,                   \\
			 & \forall \, b_j \in \calB \setminus \equi, \w_{b_j} = 0, \,
			X w = \hat y,                                                 \\
			 & \, K^\top w \leq 0, \,
			\abr{\rmin, K^\top w} = 0
			\bigg\}
		\end{aligned}
	\]
	Take any \( w' \in \calX \).
	If \( \wi' \neq 0 \), then \( \bi \in \equi \) and \( \norm{\vi}_2 = \lambda \)
	so that we may write
	\begin{align*}
		\wi' = \alpha_\bi \vi
		 & \implies \frac{\wi'}{\norm{\wi'}_2} = \frac{\alpha_\bi \vi}{\norm{\alpha_\bi \vi}_2} \\
		 & \implies \frac{\wi'}{\norm{\wi'}_2} = \frac{\vi}{\lambda}                            \\
		 & \implies \lambda \frac{\wi'}{\norm{\wi'}_2} = \vi.
	\end{align*}
	That is,
	\[
		\Xbi^\top (Xw - y) + \lambda \frac{\wi'}{\norm{\wi'}_2} + \Ki \rmin_\bi = 0,
	\]
	which is exactly stationarity of the Lagrangian.

	If \( \wi' = 0 \), then
	\[
		X w' = \hat y \implies \norm{\Xbi^\top (y - Xw') + K \rmin_\bi}_2 \leq \lambda,
	\]
	which also implies the Lagrangian is stationary.

	Since \( \abr{\rmin, K^\top w'} = 0 \) and
	\( K^\top w' \leq 0 \), i.e. \( w' \) is feasible,
	it is clear that complementary slackness must also hold.
	We conclude that \( (w', \rmin) \) are primal-dual optimal by the KKT
	conditions and the proof is complete.
\end{proof}


\uniqueCGL*
\begin{proof}
	Suppose by way of contradiction that \( w, w' \in \solfn(\lambda) \)
	such that \( w \neq w' \).
	By \cref{prop:sol-fn-cgl}, we have
	\begin{align*}
		0 = X (w - w')
		 & = \sum_{\bi \in \calS_\lambda} \Xbi (\wi - \wi')                    \\
		 & = \sum_{\bi \in \calS_\lambda} (\alpha_\bi - \alpha_\bi') \Xbi \vi,
	\end{align*}
	which implies that the vectors \( \cbr{\Xbi \vi}_{\calS_\lambda} \) are
	linearly dependent.

	Necessity follows from \cref{alg:pruning-solutions},
	which shows that, given a solution \( w \in \solfn(\lambda) \),
	linear dependence of \( \cbr{\Xbi w(\lambda) : \bi \in \calA(w(\lambda)) } \)
	implies the exist of at least one additional solution.
\end{proof}

\GGP*
\begin{proof}
	Suppose the group Lasso solution is not unique.
	Then, \cref{cor:uniqueness-cond-cgl} implies
	\[
		\calN_\lambda = \Null(\Xe) \bigcap \cbr{z : \zi = \alpha_\bi \ci, \bi \in \equi },
	\]
	is non-empty.
	That is, there exist \( \alpha_{\bi} \geq 0 \) such that
	\begin{align*}
		X_{b_j} c_{b_j}
		 & = \sum_{\bi \in \equi \setminus j} \alpha_{\bi} \Xbi \ci.
		\intertext{Taking inner-products on both sides with the residual
			\( r \),}
		r^\top X_{b_j} c_{b_j}
		 & = \sum_{\bi \in \equi \setminus j} \alpha_{\bi} r^\top \Xbi \ci \\
		\implies
		c_{b_j}^\top c_{b_j}
		 & = \sum_{\bi \in \equi \setminus j} \alpha_{\bi} \ci^\top \ci    \\
		\implies \lambda^2
		 & = \sum_{\bi \in \equi \setminus j} \alpha_{\bi} \lambda^2       \\
		\implies 1
		 & = \sum_{\bi \in \equi \setminus j} \alpha_{\bi}.
	\end{align*}
	Thus, we deduce that
	\begin{equation}\label{eq:dependence-relation}
		X_{b_j} c_{b_j}
		= \sum_{\bi \in \equi \setminus j} \beta_{\bi} \Xbi \ci,
	\end{equation}
	where \( \sum_{\bi \in \equi \setminus j} \beta_{\bi} = 1 \),
	meaning \( X_{b_j} c_{b_j} \in \text{affine}(\Xbi \zi : \bi \in \equi \setminus b_j) \)
	Now, suppose that \( |\equi| > n + 1 \).
	Then, \( \cbr{\Xbi \ci : \bi \in \equi \setminus j} \) are linearly dependent
	and, by eliminating dependent vectors \( \Xbi \ci \),
	we can repeat the above proof with a subset \( \calE' \) of at most \( n + 1 \)
	blocks.
	Noting \( \norm{\ci}_2 = \lambda \) for each \( \bi \in \equi \) and
	rescaling both sides of \cref{eq:dependence-relation} by \( \lambda \)
	implies the existence of unit vectors \( z_{\bi} \) which contradict
	GGP.
	This completes the proof.
\end{proof}

\begin{proposition}\label{prop:general-position-relation}
	Group general position does not imply
	the columns of \( X \) are in general position.
	Similarly, general position of the columns of \( X \) does not imply
	group general position.
\end{proposition}
\begin{proof}
	Let \( x_1, \ldots, x_{d} \in \R^{d-1} \).
	Consider the simple case where we have two groups: \( b_1 = \cbr{1} \)
	and \( b_2 = \cbr{2, \ldots, d} \).
	Group general position is violated if there exists a unit vector \( z_{b_2} \)
	such that
	\begin{align*}
		x_{1}
		 & = X_{b_2} z_{b_2}.   \\
		\iff x_1
		 & \in X_{b_2} B_{d-1},
	\end{align*}
	where \( B_{d-1} = \cbr{z \in \R^{d-1} : \norm{z}_2 \leq 1} \).
	In contrast, general position is violated if
	\begin{align*}
		x_1
		 & \in \text{affine}(x_2, \ldots, x_{d}) \\
		\iff x_1
		 & \in X_{b_2} \cbr{z : \abr{z, 1} = 1}.
	\end{align*}
	Taking \( X_{b_2} = I \), it is trivial to see that group general position can
	hold when general position is violated and vice-versa.
\end{proof}
